% EDITED 2026-09-12 scripts/patch_input_zd_appendix_2026-09-12T2000.py
% =====================================================================
%  APPENDIX: The zero-divisor structure of the sedenions
%  Drafted 2026-09-01. This is the tex consumer of the refs-first batch
%  zd_structure_refs_2026-08-31.{ris,bib}; all five keys are cited below:
%  moreno1998zerodivisors, biss2008annihilators, cawagas2004sedenion,
%  demarrais2000assessors, gunaydin1973quark.
%
%  PLACEMENT: appendix material. Either paste after \appendix in main.tex
%  or \input this file there. Section label minted: sec:zd-structure.
%  Labels minted (grep the tree for collisions before compiling):
%    sec:zd-structure, eq:zd-defect, prop:zd-locus, prop:zd-defect,
%    prop:zd-census, prop:zd-aut, prop:zd-colour.
%  CROSS-REFS used (verified to exist per the Higgs-subsection header):
%    sec:masses, sec:generations. No others.
%  DEPENDENCY NOTE: eq:zd-defect intentionally restates the identity that
%    the bridge splice mints as eq:debt-defect in sec:masses; an appendix
%    should be self-contained. If both land, add after eq:zd-defect:
%    "(quoted in \S\ref{sec:masses} as Eq.~\eqref{eq:debt-defect})".
%  CONVENTION: Cayley--Dickson product (a,b)(c,d) = (ac - \bar d b, da + b \bar c);
%    e4 the octonion doubling unit, e8 = ell the tag; orientation sign of
%    the coassociative form sigma = +1 in this convention.
%  VERIFICATION LEDGER (not printed): every identity below is machine-
%    verified -- sedenion_zero_divisor_scan_2026-08-31.py (echo owed),
%    sedenion_debt_koide_scan_2026-08-31.py and
%    sedenion_zd_colour_scan_2026-08-31.py (double-verified),
%    defect_calibration_scan_2026-09-01.py, accidental_u1_scan_2026-09-01.py,
%    zd_sterile_scan_2026-09-01.py (single environment, echoes owed).
%  PRIORITY CHECK OWED before submission: eq:zd-defect and the line-typing
%    of Proposition prop:zd-census are elementary and may exist in the
%    Cayley--Dickson literature (Moreno-adjacent); a pass is owed before
%    any novelty is implied. The prose below claims none.
%  British English; ASCII only; environments proposition/remark/proof
%    assumed present (amsthm), as elsewhere in the paper.
% =====================================================================

\section{The zero-divisor structure of the sedenions}
\label{sec:zd-structure}

The sedenions $\mathbb{S}$, the sixteen-dimensional Cayley--Dickson double of the octonions, enter
the main text twice: as the algebra whose three octonionic halvings carry the generations, and as
the address of the mixing seat. Both roles rest on the structure of the zero divisors, which this
appendix assembles. Throughout, $\mathbb{S} = \mathbb{O} \oplus \mathbb{O}\ell$ with the doubling
unit $\ell$ the tag; every statement has additionally been verified by explicit numerical encoding,
and the derivations are sketched in full outline.

\begin{proposition}[The locus]\label{prop:zd-locus}
The zero divisors of unit norm in $\mathbb{S}$ are exactly the elements $z = x + y\ell$ with
$x, y \in \mathrm{Im}\,\mathbb{O}$ orthonormal. Writing $j = xy$, one has $jx = y$ and $jy = -x$,
so $\mathrm{span}(x,y)$ is a complex line for left multiplication by $j$, and the annihilator is
\[
  \mathrm{Ann}(z) \;=\; \{\, u + (uj)\ell \;:\; u \perp \langle 1, x, y, j \rangle \,\},
\]
of real dimension exactly four, itself consisting of zero divisors whose direction is $-j$. The
unit zero divisors form the Stiefel manifold $V_{7,2} = G_2/\mathrm{SU}(2)$, and the set of unit
zero-divisor pairs is a $G_2$-torsor, homeomorphic to $G_2$ itself.
\end{proposition}

\begin{proof}[Proof sketch]
That every such $z$ annihilates the displayed set is a direct computation in the doubled product;
$jx = y$ follows from flexibility and alternativity. That nothing else is a zero divisor, and the
identification of the pair set with $G_2$, are due to Moreno \autocite{moreno1998zerodivisors};
the annihilator dimension is the case $n = 4$ of the bound $2^n - 4n + 4$ of Biss, Dugger and
Isaksen \autocite{biss2008annihilators}, here attained uniformly. Transitivity of $G_2$ on
orthonormal pairs of imaginary octonions gives $V_{7,2}$ with stabiliser $\mathrm{SU}(2)$, and the
$\mathrm{SU}(2)$ acts simply transitively on each annihilator sphere, giving the torsor.
\end{proof}

\begin{proposition}[Norm defect and calibration]\label{prop:zd-defect}
For all $a, b, c, d \in \mathbb{O}$,
\begin{equation}
  \bigl\lvert (a,b)(c,d) \bigr\rvert^{2}
  - \bigl\lvert (a,b) \bigr\rvert^{2} \bigl\lvert (c,d) \bigr\rvert^{2}
  \;=\; 2\,\bigl\langle a,\, [c,b,d] \bigr\rangle ,
  \label{eq:zd-defect}
\end{equation}
and the right-hand side equals, up to a fixed orientation sign, four times the coassociative
calibration four-form of the $G_2$-structure evaluated on the four components (quoted in \S\ref{sec:masses} as Eq.~\eqref{eq:debt-defect}). Consequently the
product-norm ratio $\lvert zw \rvert / (\lvert z\rvert\,\lvert w\rvert)$ ranges over
$[0, \sqrt{2}\,]$, both extremes attained exactly on configurations whose imaginary components span
a coassociative four-plane with $\lvert a\rvert = \lvert b\rvert$ and
$\lvert c\rvert = \lvert d\rvert$: the zero divisors are the calibration bound saturated in the
annihilating orientation, and $\sqrt2$ is the same bound in the amplifying one.
\end{proposition}

\begin{proof}[Proof sketch]
Expand both squared norms in the doubled product; the diagonal terms cancel against the product of
norms, and the two cross terms differ by one re-bracketing, which the trace identities
$\mathrm{Re}((xy)z) = \mathrm{Re}(x(yz))$ and $\mathrm{Re}(xy) = \mathrm{Re}(yx)$ convert into the
associator paired with $a$; real parts drop by alternation. The identification with the
coassociative form is the classical relation between the associator and the dual of
$\varphi(x,y,z) = \langle xy, z\rangle$, and the range statement is the Hadamard-type bound with
its equality case, the calibration bound of calibrated geometry.
\end{proof}

\begin{proposition}[The finite census]\label{prop:zd-census}
Index the imaginary units of $\mathbb{S}$ by the points of $\mathrm{PG}(3,2)$, so that lines are
quaternion subalgebras and planes are the fifteen eight-dimensional unit-spanned subalgebras. Then
eight planes are octonion algebras --- $\mathbb{O}$ and the seven doubles
$\mathbb{H}_L \oplus \mathbb{H}_L\ell$ --- and the remaining seven,
$Q_L = \mathbb{H}_L \oplus \mathbb{H}_L^{\perp}\ell$, are the quasi-octonion subalgebras carrying
the zero divisors \autocite{cawagas2004sedenion,demarrais2000assessors}. Sorting the thirty-five
lines by how many of their three planes are division algebras: the seven lines through the tag are
$(3,0)$; the seven Fano lines inside $\mathbb{O}$ are $(2,1)$; the twenty-one mixed lines are
$(1,2)$. In particular the only quaternion cores all three of whose extensions are division
algebras are those containing the tag, $\langle 1, j, \ell, j\ell\rangle$, and since automorphisms
preserve zero divisors, no automorphism of $\mathbb{S}$ permutes the triple of extensions of any
core inside $\mathbb{O}$: a symmetric generation triple can only be built on the tag.
\end{proposition}

\begin{proposition}[The automorphism circle and its breaking]\label{prop:zd-aut}
Let $\psi_\theta$ act as the identity on the tag plane $\langle 1, \ell\rangle$ and as right
multiplication by $\lambda = \exp(\theta\ell)$ on its fourteen-dimensional complement. For
imaginary $p, q$ and unit $\mu, \nu$ in the tag plane,
$(p\mu)(q\nu) = -\langle p, q\rangle\,\bar\mu\nu + (p \times q)\,\bar\mu\bar\nu$, so the tag phase
is conserved by the product only modulo three, and $\psi_\theta$ is an automorphism precisely when
$\lambda^3 = 1$, with automorphism defect $2\lvert \sin(3\theta/2)\rvert$ per unit cross product
otherwise. The surviving $\mathbb{Z}_3$, together with the tag reflection $\tau$, generates an
$S_3$ commuting with the diagonal $G_2$; its fixed space is exactly the tag plane. Every rotated
copy $\mathbb{O}\lambda$ is pairwise norm-multiplicative, but only the three with $\lambda^3 = 1$
close under multiplication: the halvings. The zero-divisor locus is invariant under the whole tag
circle, of which the algebra retains only the $\mathbb{Z}_3$; the halvings contain no zero
divisors, and the locus is exactly the maximally cross-halving elements.
\end{proposition}

\begin{proposition}[The colour dictionary]\label{prop:zd-colour}
Fix a unit imaginary $j$ and let $\mathrm{SU}(3) = \mathrm{Stab}_{G_2}(j)$ act on
$\mathbb{C}^3_j = j^{\perp} \subset \mathrm{Im}\,\mathbb{O}$ in the fundamental
\autocite{gunaydin1973quark}. Then the unit zero divisors with direction $j$ are exactly the
tag-doubled unit colour vectors $z = x + (jx)\ell$, an $S^5$, so the locus fibres over
$S^6 = G_2/\mathrm{SU}(3)$ with colour spheres as fibres; the tag phase on a zero divisor is the
colour phase of $x$; the annihilator of $z$ is the perpendicular anti-colour plane at $-j$; the
product of a colour and an anti-colour element lies in the core
$\langle 1, j, \ell, j\ell\rangle$ with squared norm $4\lvert h(x,u)\rvert^2$, the hermitian
overlap, the adjoint channel being absent from the algebra, while colour with colour returns the
third colour line through the antisymmetric channel. A zero-divisor pair over $j$ is an
orthonormal colour two-frame, so the fibre of Moreno's $G_2$ over $S^6$ is the colour group
itself; and the generation $\mathbb{Z}_3$, restricted to the locus, coincides with the centre of
$\mathrm{SU}(3)_j$, from which it differs off the locus by the chord $\sqrt3 = \lvert \omega - 1
\rvert$ on the direction itself.
\end{proposition}

\begin{remark}[Invertible elements, lossy channels]
Every unit zero divisor satisfies $z^2 = -1$: it is a two-sided invertible element whose left
multiplication nevertheless has a four-dimensional kernel, so that $z^{-1}(zw) = 0 \neq w$ on the
annihilator. In $\mathbb{S}$, invertibility of the element and invertibility of the channel come
apart; the loss is a property of the composition, not of the state.
\end{remark}

\begin{remark}[What this appendix does not do]
The norm defect of Eq.~\eqref{eq:zd-defect} against any fixed frame is invariant under the whole
tag circle, because the form is alternating in the pair the circle rotates; the three halvings are
isometric under the $\mathbb{Z}_3$. Nothing in this appendix therefore distinguishes the
generations intrinsically, and no functional built from the bare defect can split their masses.
The splitting is relational --- one selected frame reading three isometric copies --- and is the
business of \S\ref{sec:masses}; the identification of the halvings with the generations and of the
locus with the mixing seat is made in \S\ref{sec:generations}. Here the algebra is only permitted
to say what it can prove: where the debt settles, where it cannot, and why the triple sits on the
tag.
\end{remark}
% ================================================================ END
